% =========================================================================
% Secao 8 - Conclusoes e agenda
% =========================================================================

\section{Conclusões e agenda}
\label{sec:conclusao}

Esta nota propõe a construção de uma família de indicadores de
fluxo escolar, a saber, abandono, evasão, promoção, repetência,
migração para EJA e entrada no sistema, a partir do painel
longitudinal de cinco trimestres da PNAD Contínua. A literatura prévia brasileira sobre
fluxo escolar, em particular o trabalho do grupo PROFLUXO
\citep{Fletcher1985_modelo, Ribeiro1991_pedagogia,
KleinRibeiro1991_censo, Klein2006_educacao}, abordou o problema com
instrumentos diferentes, baseados em um modelo de fluxo de estado
estacionário aplicado a um corte transversal da PNAD. Esta nota não
é uma continuação ou extensão desse trabalho. Não há sobreposição
de autores e o arcabouço técnico é completamente distinto. O ponto
de contato é apenas o objeto comum, qual seja, usar pesquisa
domiciliar para medir fluxo escolar, e a importância histórica que
o PROFLUXO tem nas estatísticas educacionais brasileiras, em
particular a evidência apresentada por \cite{Ribeiro1991_pedagogia}
sobre a pedagogia da repetência no Brasil.

A análise mostra que a PNADC oferece quatro vantagens estruturais
sobre as estatísticas oficiais do INEP. A primeira é a
tempestividade, pois a PNADC é divulgada trimestralmente, com
defasagem máxima de quatro a cinco meses. Os Indicadores de
Trajetória do INEP, em contraste, são divulgados em ritmo errático,
com a última transição publicada, referente a 2021--2022, datada de
julho de 2025, e sem calendário público para atualizações
posteriores. Em meados de 2026, a defasagem do indicador mais
recente do INEP é de quatro anos, ao passo que a PNADC oferece
estimativas com defasagem inferior a um trimestre. A segunda
vantagem é a desagregação socioeconômica, pois a PNADC tem
variáveis individuais de renda, raça, idade, sexo e, na Visita 5,
recebimento de Bolsa Família. Os indicadores publicados pelo INEP
agregam apenas por etapa, rede, UF e localização urbana ou rural,
e os gradientes por quintil de renda e por perfil de elegibilidade
ao BFA documentados na Seção~\ref{sec:heterogeneidade} são, à data
desta nota, inviáveis com os dados publicados oficialmente. A
terceira vantagem é a independência da fonte, pois pesquisa
domiciliar e registro administrativo são metodologicamente
distintos, de modo que a convergência (ou a discrepância
controlada) entre os dois fortalece a confiança na medida e a
divergência sinaliza áreas onde o registro administrativo pode
estar viesado. Por fim, a quarta vantagem é a captura de retorno,
uma vez que a PNADC acompanha o indivíduo no domicílio, vendo
migrações entre redes, modalidades e UFs. A decomposição da
Seção~\ref{sec:inep} quantifica diretamente o componente $R$ do
\emph{gap} entre INEP e PNADC, ou seja, a fração de alunos
classificados como evadidos pelo Censo que continuam estudando
segundo a pesquisa domiciliar.

A abordagem proposta tem três limitações principais que devem
moderar a interpretação. A primeira diz respeito ao atrito do
painel. O acompanhamento longitudinal funciona apenas para
indivíduos que permanecem no mesmo domicílio, e indivíduos que
migram, por motivos como adolescentes que saem do domicílio dos
pais, casamentos ou óbitos, saem do painel. A
Seção~\ref{ssec:atrito} mostra que o atrito é correlacionado com
idade e renda, precisamente as dimensões mais relevantes para fluxo
escolar. Embora a reponderação por \emph{inverse propensity},
detalhada no Apêndice~\ref{app:variaveis}, corrija parte do viés,
esta é uma limitação inerente à pesquisa domiciliar. A segunda
limitação refere-se ao tamanho amostral em células finas. Para
subgrupos pequenos, como indígenas no EM por UF, a precisão das
estimativas é baixa, e recomendamos sempre reportar IC 95\% e
evitar agregações com $n < 100$. A terceira limitação diz respeito
ao conceito de frequência escolar. A variável \texttt{V3002} da
PNADC captura a situação na semana de referência, e em períodos de
férias ou paradas (o $Q_1$ do calendário civil é o início do ano
letivo brasileiro, e $Q_3$ e $Q_4$ abrigam férias escolares
parciais), a interpretação requer cuidado. Discutimos
especificamente esse ponto na Seção~\ref{ssec:efeito_ferias}.

A descontinuidade dos Indicadores de Trajetória do INEP cria um
vácuo informacional que esta nota busca preencher \emph{ad hoc}.
Em vista da relevância política e analítica das medidas de fluxo,
especialmente no contexto de políticas focalizadas como o
Pé-de-Meia \citep{Brasil2024_PeDeMeia}, o Bolsa Família e o
Auxílio Brasil, propomos uma sistematização institucional, isto é,
um Boletim Trimestral de Fluxo Escolar produzido com base na
PNADC, com periodicidade alinhada à divulgação dos microdados pelo
IBGE. O boletim seria, conceitualmente, o complemento contemporâneo
do PROFLUXO original, ou seja, usar pesquisa domiciliar para gerar
medidas longitudinais que o registro administrativo não consegue
manter. A diferença, agora, é que a observação substituiria o
modelo, e a frequência seria trimestral em vez de anual ou
plurianual.

A agenda subsequente inclui quatro frentes. A primeira é
aprofundar a comparação INEP \emph{versus} PNADC por UF e
macrorregião, em vez de apenas Brasil agregado. A segunda é
explorar o suplemento da quinta visita anual da PNADC, em
particular a variável \texttt{V5002A} de recebimento de Bolsa
Família, para análises com Bolsa Família observado em vez de
elegibilidade por renda, e variáveis adicionais sobre arranjo
familiar, gravidez e maternidade adolescente que ajudariam a
refinar a heterogeneidade da evasão. A terceira é estimar
efeitos do Pé-de-Meia sobre os indicadores de fluxo usando uma
estratégia de comparação antes e depois com controles em
coortes não-elegíveis, exercício que pertence a um próximo
trabalho. A quarta é replicar a abordagem para outros desfechos
educacionais, como frequência por modalidade EJA e transição
para o ensino superior. Em todos esses casos, a contribuição
metodológica desta nota, qual seja, usar o painel longitudinal
da PNADC como termômetro independente do fluxo escolar, é a
infraestrutura comum.
